\documentclass[10pt,a4paper]{amsart}
\usepackage[margin=19mm]{geometry}
\usepackage{amsmath,amssymb,mathtools}
\usepackage[scr=rsfs,cal=euler]{mathalfa}
\usepackage{xfrac}
\usepackage[unicode,hidelinks,bookmarks=false]{hyperref}
\newcommand{\ie}{i.e.\ }
\newcommand{\cf}{cf.\ }
\newcommand{\resp}{respectively\ }
\newcommand{\Subsection}{\subsection}
\providecommand{\nobreakdash}{\nobreak\hbox{-}}
\setlength{\emergencystretch}{2em}
\multlinegap0pt
\usepackage{enumerate}
\usepackage{amssymb}
\usepackage{mathtools}
\usepackage{xcolor}
\usepackage{float}
\usepackage{csquotes}
\usepackage{tikz}
\usepackage{bm}
\newtheorem{prop}{Proposition}
\newcommand{\R}{\mathbb{R}}
\title{Quadrature rules with few nodes supported on algebraic curves}
\author{Cordian Riener, Ettore Teixeira Turatti}
\date{Source: arXiv:2509.06643v3 (CC BY 4.0)}
\begin{document}
\maketitle

\noindent\emph{Source and licence.} Quadrature rules with few nodes supported on algebraic curves. Version arXiv:2509.06643v3 (CC BY 4.0), distributed by arXiv under the Creative Commons Attribution 4.0 International licence, http://creativecommons.org/licenses/by/4.0/, as recorded in the arXiv licence metadata for this exact version and shown on the abstract page https://arxiv.org/abs/2509.06643v3. Full licence notice: \texttt{licenses/arxiv-2509.06643-CC-BY-4.0.txt}.\par\smallskip

\noindent\textit{Editorial scope.} Counted unit: the article's prop prop:ttd (source0.tex lines 696--698). The article's complete original proof of it is delivered with it (source0.tex lines 699--783, one \texttt{proof} environment of 85 lines). No mathematics is written, repaired or re-derived: the statement and the proof are the article's own lines, in the article's own notation, and no retained line is substituted or re-flowed.  No further source-local context is needed: every cross-reference of every retained line resolves inside the two delivered spans.  Nothing was omitted from any retained span, so the package carries no omission record.  No retained line contains a citation, so the wrapper carries no bibliography and no bibliography entry.  No retained line contains a figure environment, an image-inclusion command or drawing code, so no image is delivered and the package declares no asset dependency.  Editorial formatting only, in addition: an AMS \texttt{amsart} preamble that restates the article's own declarations and macros used by the retained lines, namely the environments \texttt{prop} on separate counters, together with the standard packages those lines need.  The retained environments print in this document as Proposition 1 in order of appearance, whereas the article prints the same items with the numbering of its own full document.  The complete native source of the article as distributed in the arXiv e-print for this exact version is bundled unchanged as \texttt{sources/184661/source0.tex}.
\begin{prop}\label{prop:ttd}
     Let $\mathcal C\subset \R^2$ be the curve defined by $y=x^d$. {Suppose $Q(\mu)_s$ has a quadrature rule on $\mathcal C$ such that all nodes $(x_i,x_i^d)$ are such that $x_i\geq 0$}. Then for $s\geq d$, $Q(\mu)_s$ has a quadrature rule on $\mathcal C$ with at most $\left\lceil \frac{ds-1}{2}-\frac{d(d-3)}{4} \right\rceil+1$ nodes.
 \end{prop}

 \begin{proof}

The curve $\mathcal C$ is parametrized by $(t,t^d)$.  
Therefore, a quadrature rule supported on $\mathcal C$ with nodes $(t_i,t_i^d)$ satisfies
\[
\sum_{i=1}^N \omega_i x_i^a y_i^b = \sum_{i=1}^N \omega_i t_i^{a+bd},
\qquad (a+b \le s).
\]
We formulate the following optimization problem for the nodes $t_i \ge 0$ and weights $\omega_i > 0$:
\[
\begin{cases}
\min \sum_{i=1}^N t_i^2,\\[2pt]
\sum_{i=1}^N \omega_i t_i^{a+bd} = m_{a,b}, \qquad a+b \le s.
\end{cases}
\]
Here we note that  existence of such a rule follows from the hypothesis.

Again, we observe the Lagrange conditions: 
The Lagrangian is  
\[
\mathcal L(t_i,\omega_i,\lambda_{a,b})
= \sum_{i=1}^N t_i^2
+ \sum_{a+b \le s} \lambda_{a,b}\Bigl(\sum_i \omega_i t_i^{a+bd} - m_{a,b}\Bigr).
\]
Stationarity yields
\begin{align*}
0=&\frac{\partial \mathcal L}{\partial \lambda_{a,b}} = \sum_i \omega_i t_i^{a+bd}-m_{a,b},\\
0=&\frac{\partial \mathcal L}{\partial \omega_i}\phantom{a} = \sum_{a+b\le s} \lambda_{a,b}t_i^{a+bd},\\
0=&\frac{\partial \mathcal L}{\partial t_i}\phantom{a} = \omega_i\sum_{a+b\le s}(a+bd)\lambda_{a,b}t_i^{a+bd-1}+2t_i.
\end{align*}

Let  
\[
H(t) = \sum_{a+b\le s} \lambda_{a,b}\, t^{a+bd}.
\]
Then the second condition gives $H(t_i)=0$ and the third becomes
\[
H'(t_i) = -\frac{2}{\omega_i} t_i \le 0.
\]
Thus, among the positive roots of $H$, only those with non-positive derivative are candidates for local minima; at most half of the roots can satisfy this sign condition.

We again aim to bound the number of roots via dimension arguments:  Observe that the parametrization induces a linear map
\[
\psi:\R[x,y]_{\le s} \longrightarrow \R[t]_{\le sd},
\qquad p(x,y)\longmapsto p(t,t^d).
\]
Its kernel is the principal ideal $(y-x^d)$, so we have a factorization  
\[
\overline{\psi}:(\R[x,y]/(y-x^d))_{\le s} \hookrightarrow \R[t]_{\le sd}
\]
whose image is spanned by the monomials $t^{a+bd}$ with $a+b \le s$.  

The dimension formula
\[
\dim \R[x,y]_{\le s} = {s+2\choose 2}
\]
and
\[
\dim (y-x^d)\cdot \R[x,y]_{\le s-d} = {s-d+2\choose 2}
\]
gives
\[
\dim \mathrm{im}(\psi)
= {s+2\choose 2}-{s-d+2\choose 2}
= ds - \frac{d(d-3)}{2}.
\]
Hence, any $H(t)$ constructed as above has at most  
\[
ds - \frac{d(d-3)}{2} - 1
\]
positive roots by Descartes' rule of signs.
       
Among these roots, at most half satisfy $H'(t_i)\le 0$.  
Thus, the number of local minima is at most
\[
\left\lceil \frac{ds - \frac{d(d-3)}{2} - 1}{2}\right\rceil
= \left\lceil \frac{ds-1}{2}-\frac{d(d-3)}{4}\right\rceil.
\]
Including the boundary point $t=0$ from the closed interval $[0,\infty)$ adds one more node, giving
\[
N \;\le\; \left\lceil \frac{ds-1}{2}-\frac{d(d-3)}{4} \right\rceil + 1.    \]
    


 \end{proof}

\end{document}
